% ============================================================
%  PLAN DE TESIS (versión exposición) — plan_tesis_expo
%  Estructura por capítulos (plantilla UNSAAC):
%    I.  Introducción (sistema de 6 bloques)
%    II. Marco teórico y estado del arte
%    III. Metodología de la investigación
%    IV. Estudio de factibilidad
%    V.  Aspectos administrativos
%    + Matriz de consistencia (anexo)
%  Compilar: pdflatex -> biber -> pdflatex -> pdflatex
% ============================================================
\documentclass[12pt,a4paper]{article}
% ============================================================
%  Preambulo compartido - APA7 + espanol + graficos
%  Se incluye con % ============================================================
%  Preambulo compartido - APA7 + espanol + graficos
%  Se incluye con % ============================================================
%  Preambulo compartido - APA7 + espanol + graficos
%  Se incluye con \input{../preambulo.tex} en cada documento.
%  Compilar con: pdflatex -> biber -> pdflatex -> pdflatex
% ============================================================

% ---- Codificacion e idioma ----
\usepackage[utf8]{inputenc}
\usepackage[T1]{fontenc}
\usepackage{lmodern}
\usepackage[spanish,es-tabla]{babel}
\usepackage{csquotes}

% ---- Geometria (margenes tipo tesis; ajustar a SINVIC) ----
\usepackage[a4paper,left=3cm,right=2.5cm,top=2.5cm,bottom=2.5cm]{geometry}

% ---- Interlineado 1.5 ----
\usepackage{setspace}
\onehalfspacing

% ---- Tipografia de titulos ----
\usepackage{titlesec}
\titleformat{\section}{\normalfont\Large\bfseries}{\thesection}{1em}{}
\titleformat{\subsection}{\normalfont\large\bfseries}{\thesubsection}{1em}{}

% ---- Tablas y figuras ----
\usepackage{graphicx}
\usepackage{booktabs}
\usepackage{array}
\usepackage{multirow}
\usepackage{tabularx}
\usepackage{float}
\usepackage[font=small,labelfont=bf]{caption}

% ---- Graficos vectoriales (pgfplots / tikz) ----
\usepackage{pgfplots}
\pgfplotsset{compat=1.18}
\usepackage{tikz}
\usetikzlibrary{shapes.geometric,arrows.meta,positioning,calc}

% ---- Colores ----
\usepackage{xcolor}
\definecolor{azulUNSAAC}{RGB}{0,71,133}
\definecolor{grisSuave}{RGB}{240,240,240}

% ---- Bibliografia APA7 (biblatex + biber) ----
\usepackage[
  backend=biber,
  style=apa,
  sortcites=true,
  sorting=nyt,
  language=spanish
]{biblatex}
\DeclareLanguageMapping{spanish}{spanish-apa}
\addbibresource{../referencias.bib}

% ---- Listas e indices (algoritmos, ecuaciones, anexos; tablas largas) ----
% IMPORTANTE: tocloft debe cargarse ANTES de hyperref.
\usepackage{longtable}
\usepackage{algorithm}
\usepackage{algpseudocode}
\usepackage{tocloft}

% ---- Hipervinculos ----
\usepackage[hidelinks]{hyperref}
\hypersetup{
  colorlinks=true,
  linkcolor=black,
  citecolor=azulUNSAAC,
  urlcolor=azulUNSAAC
}

% ---- Icono ORCID (dibujado con tikz, compatible con pdflatex) ----
\usepackage{orcidlink}

% ---- Encabezado / pie de pagina ----
\usepackage{fancyhdr}
\pagestyle{fancy}
\fancyhf{}
\fancyfoot[C]{\thepage}
\renewcommand{\headrulewidth}{0pt}

% ---- Sangria de primer parrafo ----
\usepackage{indentfirst}
 en cada documento.
%  Compilar con: pdflatex -> biber -> pdflatex -> pdflatex
% ============================================================

% ---- Codificacion e idioma ----
\usepackage[utf8]{inputenc}
\usepackage[T1]{fontenc}
\usepackage{lmodern}
\usepackage[spanish,es-tabla]{babel}
\usepackage{csquotes}

% ---- Geometria (margenes tipo tesis; ajustar a SINVIC) ----
\usepackage[a4paper,left=3cm,right=2.5cm,top=2.5cm,bottom=2.5cm]{geometry}

% ---- Interlineado 1.5 ----
\usepackage{setspace}
\onehalfspacing

% ---- Tipografia de titulos ----
\usepackage{titlesec}
\titleformat{\section}{\normalfont\Large\bfseries}{\thesection}{1em}{}
\titleformat{\subsection}{\normalfont\large\bfseries}{\thesubsection}{1em}{}

% ---- Tablas y figuras ----
\usepackage{graphicx}
\usepackage{booktabs}
\usepackage{array}
\usepackage{multirow}
\usepackage{tabularx}
\usepackage{float}
\usepackage[font=small,labelfont=bf]{caption}

% ---- Graficos vectoriales (pgfplots / tikz) ----
\usepackage{pgfplots}
\pgfplotsset{compat=1.18}
\usepackage{tikz}
\usetikzlibrary{shapes.geometric,arrows.meta,positioning,calc}

% ---- Colores ----
\usepackage{xcolor}
\definecolor{azulUNSAAC}{RGB}{0,71,133}
\definecolor{grisSuave}{RGB}{240,240,240}

% ---- Bibliografia APA7 (biblatex + biber) ----
\usepackage[
  backend=biber,
  style=apa,
  sortcites=true,
  sorting=nyt,
  language=spanish
]{biblatex}
\DeclareLanguageMapping{spanish}{spanish-apa}
\addbibresource{../referencias.bib}

% ---- Listas e indices (algoritmos, ecuaciones, anexos; tablas largas) ----
% IMPORTANTE: tocloft debe cargarse ANTES de hyperref.
\usepackage{longtable}
\usepackage{algorithm}
\usepackage{algpseudocode}
\usepackage{tocloft}

% ---- Hipervinculos ----
\usepackage[hidelinks]{hyperref}
\hypersetup{
  colorlinks=true,
  linkcolor=black,
  citecolor=azulUNSAAC,
  urlcolor=azulUNSAAC
}

% ---- Icono ORCID (dibujado con tikz, compatible con pdflatex) ----
\usepackage{orcidlink}

% ---- Encabezado / pie de pagina ----
\usepackage{fancyhdr}
\pagestyle{fancy}
\fancyhf{}
\fancyfoot[C]{\thepage}
\renewcommand{\headrulewidth}{0pt}

% ---- Sangria de primer parrafo ----
\usepackage{indentfirst}
 en cada documento.
%  Compilar con: pdflatex -> biber -> pdflatex -> pdflatex
% ============================================================

% ---- Codificacion e idioma ----
\usepackage[utf8]{inputenc}
\usepackage[T1]{fontenc}
\usepackage{lmodern}
\usepackage[spanish,es-tabla]{babel}
\usepackage{csquotes}

% ---- Geometria (margenes tipo tesis; ajustar a SINVIC) ----
\usepackage[a4paper,left=3cm,right=2.5cm,top=2.5cm,bottom=2.5cm]{geometry}

% ---- Interlineado 1.5 ----
\usepackage{setspace}
\onehalfspacing

% ---- Tipografia de titulos ----
\usepackage{titlesec}
\titleformat{\section}{\normalfont\Large\bfseries}{\thesection}{1em}{}
\titleformat{\subsection}{\normalfont\large\bfseries}{\thesubsection}{1em}{}

% ---- Tablas y figuras ----
\usepackage{graphicx}
\usepackage{booktabs}
\usepackage{array}
\usepackage{multirow}
\usepackage{tabularx}
\usepackage{float}
\usepackage[font=small,labelfont=bf]{caption}

% ---- Graficos vectoriales (pgfplots / tikz) ----
\usepackage{pgfplots}
\pgfplotsset{compat=1.18}
\usepackage{tikz}
\usetikzlibrary{shapes.geometric,arrows.meta,positioning,calc}

% ---- Colores ----
\usepackage{xcolor}
\definecolor{azulUNSAAC}{RGB}{0,71,133}
\definecolor{grisSuave}{RGB}{240,240,240}

% ---- Bibliografia APA7 (biblatex + biber) ----
\usepackage[
  backend=biber,
  style=apa,
  sortcites=true,
  sorting=nyt,
  language=spanish
]{biblatex}
\DeclareLanguageMapping{spanish}{spanish-apa}
\addbibresource{../referencias.bib}

% ---- Listas e indices (algoritmos, ecuaciones, anexos; tablas largas) ----
% IMPORTANTE: tocloft debe cargarse ANTES de hyperref.
\usepackage{longtable}
\usepackage{algorithm}
\usepackage{algpseudocode}
\usepackage{tocloft}

% ---- Hipervinculos ----
\usepackage[hidelinks]{hyperref}
\hypersetup{
  colorlinks=true,
  linkcolor=black,
  citecolor=azulUNSAAC,
  urlcolor=azulUNSAAC
}

% ---- Icono ORCID (dibujado con tikz, compatible con pdflatex) ----
\usepackage{orcidlink}

% ---- Encabezado / pie de pagina ----
\usepackage{fancyhdr}
\pagestyle{fancy}
\fancyhf{}
\fancyfoot[C]{\thepage}
\renewcommand{\headrulewidth}{0pt}

% ---- Sangria de primer parrafo ----
\usepackage{indentfirst}


% ----- Color de celdas / rampas -----
\usepackage{colortbl}

% ----- Colores de la "ecuación universal" del título (azul/ámbar/rojo) -----
\definecolor{colProceso}{RGB}{0,71,133}    % azul  -> Proceso / Acción
\definecolor{colMetodo}{RGB}{181,128,0}    % ámbar -> Tecnología / Método
\definecolor{colContexto}{RGB}{176,42,42}  % rojo  -> Contexto / Aplicación

% ----- Rampa estándar de los gráficos: azul claro -> navy -> turquesa -> verde
%        (mismo degradado del cronograma basado en el proyecto MusicVA) -----
\definecolor{g1}{HTML}{8AC4E6}
\definecolor{g2}{HTML}{66AEDA}
\definecolor{g3}{HTML}{4294CC}
\definecolor{g4}{HTML}{2E6DA0}
\definecolor{g5}{HTML}{1E4870}
\definecolor{g6}{HTML}{1F6476}
\definecolor{g7}{HTML}{229184}
\definecolor{g8}{HTML}{27A597}
\definecolor{g9}{HTML}{40B183}
\definecolor{g10}{HTML}{59BD6F}
\definecolor{g11}{HTML}{6AC362}
\definecolor{g12}{HTML}{7BC95A}

% ----- Numeración por capítulos (I, II, ...) y subsecciones (2.1, 2.1.1) -----
\renewcommand{\thesection}{\Roman{section}}
\renewcommand{\thesubsection}{\arabic{section}.\arabic{subsection}}
\renewcommand{\thesubsubsection}{\arabic{section}.\arabic{subsection}.\arabic{subsubsection}}
\setcounter{secnumdepth}{3}
\setcounter{tocdepth}{3}

% ----- Nombres en español de los índices (como en el plan anterior) -----
\addto\captionsspanish{\renewcommand{\contentsname}{\'Indice general}}
\floatname{algorithm}{Algoritmo}
\renewcommand{\listalgorithmname}{\'Indice de algoritmos}

% ----- Índice de ecuaciones (lista propia vía tocloft) -----
\newlistof{ecuaciones}{loe}{\'Indice de ecuaciones}
\newcommand{\regecuacion}[1]{%
  \addcontentsline{loe}{ecuaciones}{\protect\numberline{\theequation}#1}}
% Formato de las entradas del índice de ecuaciones igual al de cuadros
\let\cftecuacionesfont\cfttabfont
\let\cftecuacionespagefont\cfttabpagefont
\setlength{\cftecuacionesindent}{\cfttabindent}
\setlength{\cftecuacionesnumwidth}{\cfttabnumwidth}

% ----- Fuentes de la revisión (para los antecedentes) -----
\addbibresource{antecedentes.bib}

\begin{document}

% ================= CARÁTULA (plantilla UNSAAC) =================
\begin{titlepage}
\centering
\singlespacing

{\large UNIVERSIDAD NACIONAL DE SAN ANTONIO ABAD\\ DEL CUSCO\par}
\vspace{0.2cm}
{FACULTAD DE INGENIER{\'I}A EL{\'E}CTRICA, ELECTR{\'O}NICA, INFORM{\'A}TICA Y
M{\'E}CANICA\par}
{ESCUELA PROFESIONAL DE INGENIER{\'I}A INFORM{\'A}TICA Y DE SISTEMAS\par}
\vspace{0.6cm}

\IfFileExists{../logo_unsaac.png}{\includegraphics[width=4.6cm]{../logo_unsaac.png}}{%
  \IfFileExists{logo_unsaac.png}{\includegraphics[width=4.6cm]{logo_unsaac.png}}{}}
\vspace{0.6cm}

{\scshape Plan de Tesis\par}
\vspace{0.5cm}
\rule{\linewidth}{0.4pt}\par\vspace{0.3cm}
{\large\bfseries%
``\textcolor{colProceso}{SISTEMA BIDIRECCIONAL DE TRADUCCI{\'O}N DE VOZ
QUECHUA--ESPA{\~N}OL}
\textcolor{colMetodo}{CON IA PREENTRENADA}
\textcolor{colContexto}{PARA CONSULTAS M{\'E}DICAS EN CUSCO}''\par}
\vspace{0.3cm}\rule{\linewidth}{0.4pt}\par

\vspace{1.1cm}
\begin{flushleft}
\hspace*{0.28\textwidth}%
\begin{minipage}{0.66\textwidth}
\raggedright
{\small Para optar al t{\'i}tulo profesional de:}\\[3pt]
\hspace*{1.4em}INGENIERO INFORM{\'A}TICO Y DE SISTEMAS\\[13pt]
{\small Presentado por:}\\[3pt]
\hspace*{1.4em}Alain Armando Condori Contreras\\[13pt]
{\small C{\'o}digo de alumno:}\\[3pt]
\hspace*{1.4em}021501\\[13pt]
{\small Asesor:}\\[3pt]
\hspace*{1.4em}Mtro. Ing. Ra{\'u}l Huillca Huallparimachi
\end{minipage}
\end{flushleft}

\vfill
{\bfseries Per{\'u}, Octubre de 2026\par}
\end{titlepage}

% ================= RESUMEN =================
\section*{Resumen}
\addcontentsline{toc}{section}{Resumen}
\singlespacing\noindent
En la regi{\'o}n Cusco, m{\'a}s del 55\,\% de la poblaci{\'o}n tiene el quechua como lengua
materna y cerca de 70\,000 asegurados de EsSalud son quechuahablantes \parencite{essalud2023};
esta barrera idiom{\'a}tica afecta la calidad de la atenci{\'o}n en salud \parencite{begazo2026}.
El presente trabajo propone desarrollar y evaluar un sistema bidireccional de traducci{\'o}n de
voz Quechua--Espa{\~n}ol para consultas m{\'e}dicas, que encadena reconocimiento de voz (MMS),
traducci{\'o}n autom{\'a}tica (NLLB adaptado al quechua de Cusco, \texttt{quz}) y s{\'i}ntesis de
voz (MMS-TTS) mediante el ajuste (\emph{fine-tuning}) de modelos preentrenados de c{\'o}digo
abierto. La salida por voz en ambos sentidos busca incluir a la poblaci{\'o}n quechuahablante con
baja o nula alfabetizaci{\'o}n. La validaci{\'o}n se realizar{\'a} con consultas simuladas, sin datos
reales de pacientes, midiendo precisi{\'o}n (WER, BLEU, chrF), calidad de voz (MOS) y percepci{\'o}n
de utilidad. Se espera aportar el primer corpus cl{\'i}nico Quechua--Espa{\~n}ol y evidencia
emp{\'i}rica sobre la viabilidad de estas tecnolog{\'i}as en el dominio sanitario andino.

\vspace{0.3cm}
\noindent\textbf{Palabras clave:} traducci{\'o}n autom{\'a}tica, quechua, PLN, inclusi{\'o}n digital,
comunicaci{\'o}n m{\'e}dico-paciente, ASR, TTS.
\onehalfspacing

\newpage
\tableofcontents                                   % Índice general

\clearpage
\addcontentsline{toc}{section}{\'Indice de figuras}
\listoffigures                                     % Índice de figuras

\clearpage
\addcontentsline{toc}{section}{\'Indice de cuadros}
\listoftables                                      % Índice de cuadros

\clearpage
\addcontentsline{toc}{section}{\'Indice de algoritmos}
\listofalgorithms                                  % Índice de algoritmos

\clearpage                                          % Índice de símbolos
\section*{\'Indice de s{\'i}mbolos}
\addcontentsline{toc}{section}{\'Indice de s{\'i}mbolos}
\begin{longtable}{@{}p{2.6cm} p{9cm} p{2.4cm}@{}}
\toprule
\textbf{S{\'i}mbolo} & \textbf{Descripci{\'o}n} & \textbf{Unidad} \\
\midrule \endhead
$\alpha$ & Nivel de significancia estad{\'i}stica & --- \\
$p$      & Valor de probabilidad (p-valor) & --- \\
$N$      & Tama{\~n}o de muestra / n.\ de palabras de referencia & --- \\
$\beta$  & Factor de ponderaci{\'o}n recall/precisi{\'o}n en chrF & --- \\
\bottomrule
\end{longtable}

\clearpage                                          % Índice de abreviaturas
\section*{\'Indice de abreviaturas}
\addcontentsline{toc}{section}{\'Indice de abreviaturas}
\begin{longtable}{@{}p{3cm} p{12cm}@{}}
\toprule
\textbf{Sigla} & \textbf{Significado} \\
\midrule \endhead
IA      & Inteligencia artificial \\
PLN     & Procesamiento del lenguaje natural \\
ASR     & Reconocimiento autom{\'a}tico del habla (\emph{Automatic Speech Recognition}) \\
STT     & Voz a texto (\emph{Speech-to-Text}) \\
TTS     & S{\'i}ntesis de voz (\emph{Text-to-Speech}) \\
MT      & Traducci{\'o}n autom{\'a}tica (\emph{Machine Translation}) \\
NMT     & Traducci{\'o}n autom{\'a}tica neuronal \\
MMS     & \emph{Massively Multilingual Speech} (Meta) \\
NLLB    & \emph{No Language Left Behind} (Meta) \\
WER     & Tasa de error de palabra (\emph{Word Error Rate}) \\
BLEU    & \emph{Bilingual Evaluation Understudy} \\
chrF    & \emph{Character n-gram F-score} \\
MOS     & \emph{Mean Opinion Score} \\
quz     & Quechua Collao / Cusco (c{\'o}digo ISO 639-3) \\
quy     & Quechua Ayacucho (c{\'o}digo ISO 639-3) \\
INEI    & Instituto Nacional de Estad{\'i}stica e Inform{\'a}tica \\
EsSalud & Seguro Social de Salud del Per{\'u} \\
UNSAAC  & Universidad Nacional de San Antonio Abad del Cusco \\
\bottomrule
\end{longtable}

\clearpage                                          % Índice de ecuaciones
\section*{\'Indice de ecuaciones}
\addcontentsline{toc}{section}{\'Indice de ecuaciones}
\makeatletter\@starttoc{loe}\makeatother

\clearpage                                          % Glosario de términos
\section*{Glosario de t{\'e}rminos}
\addcontentsline{toc}{section}{Glosario de t{\'e}rminos}
\begin{longtable}{@{}p{3.6cm} p{11.4cm}@{}}
\toprule
\textbf{T{\'e}rmino} & \textbf{Definici{\'o}n} \\
\midrule \endhead
Corpus                  & Conjunto estructurado de textos o grabaciones usado para entrenar o evaluar modelos. \\
Fine-tuning             & Ajuste de un modelo preentrenado con datos de un dominio espec{\'i}fico. \\
Lengua de bajo recurso  & Idioma con escasez de datos digitales para entrenar modelos de PLN. \\
Modelo preentrenado     & Modelo entrenado sobre grandes vol{\'u}menes de datos, reutilizable mediante ajuste. \\
Lengua polisint{\'e}tica & Lengua que concentra mucha informaci{\'o}n gramatical en una sola palabra, como el quechua. \\
Pipeline                & Secuencia encadenada de etapas (voz $\to$ texto $\to$ traducci{\'o}n $\to$ voz). \\
\bottomrule
\end{longtable}

\clearpage
\setcounter{table}{0}\setcounter{figure}{0}

% ##################################################################
% CAPÍTULO I: INTRODUCCIÓN (sistema de 6 bloques)
% ##################################################################
\section{Introducci{\'o}n}
La introducci{\'o}n de este plan sigue el \emph{sistema de seis bloques}: parte del contexto y el
problema general, presenta la tecnolog{\'i}a disponible, identifica el vac{\'i}o de investigaci{\'o}n,
precisa el problema espec{\'i}fico, formula la propuesta con sus objetivos, hip{\'o}tesis y
justificaci{\'o}n, y cierra con la organizaci{\'o}n del documento.

\subsection{Contexto y problema general}
El Per{\'u} es un pa{\'i}s plurilingüe en el que el quechua es la lengua originaria con mayor
n{\'u}mero de hablantes. Los Censos Nacionales 2017 registran a millones de peruanos que declaran
el quechua como lengua materna, con una marcada concentraci{\'o}n en la sierra sur
\parencite{inei2017}. En la regi{\'o}n Cusco esta presencia es particularmente intensa: m{\'a}s del
55\,\% de la poblaci{\'o}n tiene el quechua como lengua materna \parencite{essalud2023}. Esta
realidad ling{\"u}{\'i}stica condiciona el acceso efectivo a los servicios p{\'u}blicos
esenciales ---la atenci{\'o}n en salud entre ellos--- cuando operan de manera predominante en
espa{\~n}ol.

En el {\'a}mbito sanitario, la distancia idiom{\'a}tica entre un personal de salud mayoritariamente
castellanohablante y una poblaci{\'o}n usuaria quechuahablante origina una barrera comunicativa con
consecuencias cl{\'i}nicas directas. EsSalud reconoce que cerca de 70\,000 de sus asegurados en
Cusco son quechuahablantes y reactiv{\'o} en 2023 la central telef{\'o}nica en quechua
\parencite{essalud2023}; sin embargo, esa medida no alcanza la interacci{\'o}n cara a cara del
consultorio, donde se recogen s{\'i}ntomas, se formula el diagn{\'o}stico y se entregan las
indicaciones. La evidencia confirma la gravedad del problema: \textcite{begazo2026} documenta
c{\'o}mo la desconexi{\'o}n idiom{\'a}tica vulnera la relaci{\'o}n terap{\'e}utica en hospitales de
Cusco, y \textcite{cuaila2022} halla una relaci{\'o}n entre el conocimiento del quechua del personal
y la calidad de atenci{\'o}n percibida. A ello se suma una barrera de alfabetizaci{\'o}n: cerca del
22.8\,\% de las mujeres rurales no sabe leer ni escribir \parencite{inei2023}, lo que obliga a una
soluci{\'o}n basada en \emph{voz} y no solo en texto (Figura~\ref{fig:barreras}).

\begin{figure}[H]
\centering
\begin{tikzpicture}
\begin{axis}[
    ybar, width=12cm, height=6cm, bar width=28pt,
    ymin=0, ymax=100, ylabel={Porcentaje (\%)},
    symbolic x coords={Lengua materna,Asegurados EsSalud,Analfab. mujer rural},
    xtick={Lengua materna,Asegurados EsSalud,Analfab. mujer rural},
    nodes near coords, nodes near coords style={font=\small},
    x tick label style={align=center,font=\small,text width=2.5cm},
    enlarge x limits=0.25, ymajorgrids=true, grid style=dashed,
]
\addplot[fill=g4, draw=g5, bar shift=0pt] coordinates {(Lengua materna,55)};
\addplot[fill=g8, draw=g7, bar shift=0pt] coordinates {(Asegurados EsSalud,15)};
\addplot[fill=g12, draw=g11, bar shift=0pt] coordinates {(Analfab. mujer rural,22.8)};
\end{axis}
\end{tikzpicture}
\caption{Magnitud de la problem{\'a}tica en Cusco. Elaboraci{\'o}n propia con datos de \textcite{essalud2023} e \textcite{inei2023}.}
\label{fig:barreras}
\end{figure}

\subsection{Tecnolog{\'i}a o m{\'e}todo}
El estado del arte en procesamiento del lenguaje natural ofrece hoy una oportunidad que no
exist{\'i}a hace pocos a{\~n}os. La estrategia dominante para lenguas de bajo recurso ya no consiste
en entrenar modelos desde cero, sino en adaptar (\emph{fine-tuning}) grandes modelos multiling{\"u}es
preentrenados. En la voz, el proyecto MMS \parencite{pratap2024mms} cubre el quechua de Cusco
(\texttt{quz}) en reconocimiento y s{\'i}ntesis; Whisper \parencite{radford2023whisper} declara una
cobertura marginal del quechua y se reserva como l{\'i}nea base. En traducci{\'o}n, NLLB
\parencite{costajussa2022nllb} incluye el quechua de Ayacucho (\texttt{quy}) y no el de Cusco, por lo
que el par debe adaptarse a \texttt{quz} ---esa brecha dialectal es el n{\'u}cleo del aporte
t{\'e}cnico---. Existen adem{\'a}s antecedentes de traducci{\'o}n neural para esta lengua
polisint{\'e}tica \parencite{ortega2020quechua,mager2021americasnlp} y recursos como el corpus
\emph{Huqariq} \parencite{zevallos2022huqariq}.

\subsection{Vac{\'i}o de investigaci{\'o}n}
Persiste una desconexi{\'o}n entre ambos mundos: el problema social est{\'a} bien documentado y la
tecnolog{\'i}a base existe, pero nadie ha unido los dos extremos. No se dispone de un corpus
cl{\'i}nico Quechua--Espa{\~n}ol ni de una evaluaci{\'o}n emp{\'i}rica del desempe{\~n}o de estos
modelos sobre el vocabulario propio de la consulta m{\'e}dica (triaje, s{\'i}ntomas e indicaciones),
y NLLB a{\'u}n no cubre la variante de Cusco (\texttt{quz}). Cerrar esa intersecci{\'o}n es el aporte
de este proyecto.

\subsection{Problema espec{\'i}fico}
El problema espec{\'i}fico de esta investigaci{\'o}n se concentra en la comunicaci{\'o}n
m{\'e}dico-paciente mediada por un sistema bidireccional de traducci{\'o}n de voz Quechua--Espa{\~n}ol
basado en modelos preentrenados, y en la medida en que su ajuste al dominio cl{\'i}nico mejora la
precisi{\'o}n y la comprensi{\'o}n frente a la ausencia de apoyo tecnol{\'o}gico.

\subsubsection{Problema general}
{\textquestiondown}De qu{\'e} manera un sistema bidireccional de traducci{\'o}n de voz
Quechua--Espa{\~n}ol, basado en modelos de IA preentrenados, influye en la comunicaci{\'o}n entre el
personal de salud y los pacientes quechuahablantes durante las consultas m{\'e}dicas en Cusco?

\subsubsection{Problemas espec{\'i}ficos}
\begin{enumerate}
  \item {\textquestiondown}Qu{\'e} caracter{\'i}sticas debe reunir un corpus de vocabulario cl{\'i}nico Quechua--Espa{\~n}ol para el ajuste de los modelos de traducci{\'o}n por voz?
  \item {\textquestiondown}En qu{\'e} medida el ajuste (\emph{fine-tuning}) de modelos preentrenados mejora la precisi{\'o}n sobre el vocabulario cl{\'i}nico?
  \item {\textquestiondown}Qu{\'e} modalidad de salida (texto para el personal; texto y voz para el paciente) resulta m{\'a}s confiable y comprensible?
  \item {\textquestiondown}Qu{\'e} nivel de utilidad, confianza y comprensi{\'o}n perciben el personal y los pacientes al usar el sistema en consultas simuladas?
\end{enumerate}

\subsection{Propuesta, objetivo y justificaci{\'o}n}

\subsubsection{Propuesta}
Se propone desarrollar y evaluar un sistema bidireccional de traducci{\'o}n de voz
Quechua--Espa{\~n}ol para consultas m{\'e}dicas, que encadena reconocimiento de voz (MMS),
traducci{\'o}n autom{\'a}tica (NLLB adaptado de \texttt{quy} a \texttt{quz}) y s{\'i}ntesis de voz
(MMS-TTS) mediante el ajuste de modelos preentrenados de c{\'o}digo abierto, con salida por voz en
ambos sentidos del di{\'a}logo.

\subsubsection{Objetivo general}
Desarrollar y evaluar un sistema bidireccional de traducci{\'o}n de voz Quechua--Espa{\~n}ol, basado
en modelos de IA preentrenados, para mejorar la comunicaci{\'o}n m{\'e}dico-paciente durante las
consultas m{\'e}dicas en Cusco.

\subsubsection{Objetivos espec{\'i}ficos}
\begin{enumerate}
  \item Construir y validar un corpus de vocabulario y frases de consulta m{\'e}dica en Quechua--Espa{\~n}ol.
  \item Adaptar (\emph{fine-tuning}) los modelos preentrenados de voz (MMS) y, como aporte central, el modelo de traducci{\'o}n NLLB del quechua de Ayacucho (\texttt{quy}) al de Cusco (\texttt{quz}) sobre el dominio cl{\'i}nico.
  \item Integrar el \emph{pipeline} bidireccional con salida por voz y texto para ambos interlocutores, utilizable por pacientes con baja o nula alfabetizaci{\'o}n.
  \item Validar el sistema en consultas simuladas, midiendo precisi{\'o}n, calidad de voz, tiempo y percepci{\'o}n.
\end{enumerate}

\subsubsection{Hip{\'o}tesis}
\textbf{General:} un sistema bidireccional de traducci{\'o}n de voz Quechua--Espa{\~n}ol, basado en
modelos de IA preentrenados, mejora significativamente la comunicaci{\'o}n m{\'e}dico-paciente durante
las consultas m{\'e}dicas en Cusco frente a la ausencia de apoyo tecnol{\'o}gico.
\begin{itemize}
  \item \textbf{H1:} un corpus cl{\'i}nico Quechua--Espa{\~n}ol validado proporciona la base l{\'e}xica necesaria para el ajuste de los modelos.
  \item \textbf{H2:} el ajuste (\emph{fine-tuning}) de los modelos preentrenados reduce el WER y mejora el BLEU/chrF frente a los modelos sin ajustar.
  \item \textbf{H3:} la salida texto+voz mejora la comprensi{\'o}n de las indicaciones frente a solo texto, sobre todo con baja alfabetizaci{\'o}n.
  \item \textbf{H4:} el uso del sistema mejora significativamente la comprensi{\'o}n mutua percibida frente a la ausencia de apoyo tecnol{\'o}gico.
\end{itemize}

\subsubsection{Justificaci{\'o}n}
\textbf{Te{\'o}rica:} aporta evidencia emp{\'i}rica del desempe{\~n}o de modelos preentrenados en una
lengua de bajo recurso y morfolog{\'i}a compleja. \textbf{Pr{\'a}ctica:} entrega un prototipo
piloteable en un centro de salud. \textbf{Social:} ataca una inequidad documentada en el acceso a la
salud \parencite{begazo2026}, en l{\'i}nea con el derecho de los pacientes a ser atendidos en su
lengua materna \parencite{essalud2023}. \textbf{Metodol{\'o}gica:} establece un proceso replicable
(corpus + ajuste + evaluaci{\'o}n) en otras variantes del quechua o en otros dominios. Entre las
contribuciones originales esperadas destacan el \textbf{primer corpus de dominio m{\'e}dico}
Quechua--Espa{\~n}ol para tecnolog{\'i}as del habla, la \textbf{adaptaci{\'o}n de NLLB de
\texttt{quy} a \texttt{quz}} en el dominio cl{\'i}nico (aporte t{\'e}cnico central) y un
\textbf{sistema integrado ASR+MT+TTS} con salida por voz para usuarios con baja alfabetizaci{\'o}n.
En conjunto, el proyecto contribuye a la equidad en el acceso a la salud de la poblaci{\'o}n
quechuahablante de Cusco, a la seguridad del acto m{\'e}dico y a la revitalizaci{\'o}n y el prestigio
digital del quechua.

\subsubsection{Alcances y limitaciones}
\noindent\textbf{Alcances.}
\begin{itemize}
  \item Par ling{\"u}{\'i}stico Quechua de Cusco (\texttt{quz}) -- Espa{\~n}ol, en ambos sentidos.
  \item Vocabulario acotado de consulta general (triaje, s{\'i}ntomas, indicaciones).
  \item Prototipo integrado ASR+MT+TTS con salida por voz y texto.
  \item Validaci{\'o}n con consultas simuladas en un centro de salud de referencia en Cusco.
\end{itemize}
\noindent\textbf{Limitaciones.}
\begin{itemize}
  \item \textbf{Temporal:} ejecuci{\'o}n entre octubre 2026 y enero 2027.
  \item \textbf{{\'E}tica:} se emplean di{\'a}logos simulados/actuados, sin datos reales de pacientes.
  \item Desempe{\~n}o condicionado por la escasez de datos de dominio cl{\'i}nico en quechua.
  \item Resultados circunscritos a la variante de Cusco; no generalizables sin ajuste a otras variantes.
\end{itemize}

\subsection{Organizaci{\'o}n de la tesis}
El documento se organiza en cinco cap{\'i}tulos. El \textbf{Cap{\'i}tulo~I} presenta la
introducci{\'o}n ---contexto, tecnolog{\'i}a, vac{\'i}o, problema, propuesta, objetivos, hip{\'o}tesis
y justificaci{\'o}n---. El \textbf{Cap{\'i}tulo~II} desarrolla el marco te{\'o}rico y el estado del
arte derivado de la revisi{\'o}n sistem{\'a}tica. El \textbf{Cap{\'i}tulo~III} detalla la
metodolog{\'i}a: enfoque, dise{\~n}o, variables, unidades de an{\'a}lisis, t{\'e}cnicas,
procedimiento, protocolo de evaluaci{\'o}n y aspectos {\'e}ticos. El \textbf{Cap{\'i}tulo~IV}
analiza la factibilidad t{\'e}cnica, econ{\'o}mica y operativa. El \textbf{Cap{\'i}tulo~V} expone los
aspectos administrativos ---cronograma, presupuesto y equipo---. Finalmente, la \textbf{matriz de
consistencia} articula el problema, los objetivos, las hip{\'o}tesis y las variables.

% ##################################################################
% CAPÍTULO II: MARCO TEÓRICO Y ESTADO DEL ARTE
% ##################################################################
\section{Marco te{\'o}rico y estado del arte}

\subsection{Antecedentes}
Los trabajos previos sobre tecnolog{\'i}as del habla y traducci{\'o}n autom{\'a}tica para el quechua
se organizan en tres niveles ---internacional, nacional y local---. En conjunto evidencian avances
s{\'o}lidos por componente (ASR, MT y TTS), pero ninguno aborda su integraci{\'o}n en el dominio
m{\'e}dico-paciente ni la variante de Cusco (\texttt{quz}).

\subsubsection{Antecedentes internacionales}
La traducci{\'o}n neural para el quechua fue abordada por primera vez por \textcite{ortega2020quechua},
que introdujo t{\'e}cnicas de segmentaci{\'o}n morfol{\'o}gica para esta lengua polisint{\'e}tica. Las
competencias AmericasNLP impulsaron el {\'a}rea: \textcite{Gow-Smith2023192} extendieron NLLB-200 con
una mejora promedio de 11\,\% en chrF, mientras que \textcite{Ahmed2023156} y \textcite{Dhawan2026119}
exploraron modelos preentrenados y datos sint{\'e}ticos \parencite{Ebrahimi2024236}. En el habla,
\textcite{Chen20246444} evaluaron representaciones autosupervisadas para lenguas ind{\'i}genas y
\textcite{Chen202120} propusieron una segmentaci{\'o}n guiada por la morfolog{\'i}a. La serie
\emph{QUESPA} en IWSLT llev{\'o} la traducci{\'o}n de voz Quechua--Espa{\~n}ol de cerca de 16 a
\textbf{26.7 BLEU} entre 2023 y 2025 \parencite{Ortega2023261,Ortega2024173,Ortega2025260}. Estos
avances se apoyan en grandes modelos multiling{\"u}es preentrenados ---MMS \parencite{pratap2024mms},
NLLB \parencite{costajussa2022nllb} y Whisper \parencite{radford2023whisper}---.

\subsubsection{Antecedentes nacionales}
En el Per{\'u}, \textcite{Zevallos202061} implementaron un ASR de quechua con HMM/HTK alcanzando
\textbf{12.70\,\% de WER}, mejorado luego mediante aumento de datos hasta 22.75\,\% de WER
\parencite{Zevallos20223518}; el mismo grupo aport{\'o} el modelo \emph{QuBERT}
\parencite{Zevallos20221} y el corpus \emph{Huqariq} \parencite{zevallos2022huqariq}. En s{\'i}ntesis
de voz, \textcite{HuamanCuro2026} evaluaron cuatro arquitecturas TTS para el quechua chanca con un MOS
de hasta \textbf{4.9}. En traducci{\'o}n, \textcite{Barbaran2026464} (UTP) y \textcite{CuevaMedina2024}
(UCSM, Arequipa) ajustaron modelos neuronales Quechua--Espa{\~n}ol, y \textcite{Veliz2021241} (PUCP) y
\textcite{Oncevay2021194} trabajaron enfoques multiling{\"u}es para las lenguas del Per{\'u}. En
traducci{\'o}n de voz, \textcite{Alvarez-Crespo2023837} (UPC) desarrollaron un subtitulado en tiempo
real Espa{\~n}ol$\to$Quechua con arquitectura en cascada, y \textcite{9514251} construyeron una
plataforma para recolectar un corpus de voz Espa{\~n}ol/Quechua mediante la API
\emph{Siminchikkunarayku}. Ninguno de estos trabajos abord{\'o} el dominio sanitario.

\subsubsection{Antecedentes locales}
En Cusco, la evidencia se concentra en el problema m{\'a}s que en la soluci{\'o}n:
\textcite{begazo2026} document{\'o} c{\'o}mo la barrera idiom{\'a}tica vulnera la relaci{\'o}n
terap{\'e}utica en hospitales de referencia de la ciudad, y EsSalud reactiv{\'o} en 2023 una central
telef{\'o}nica en quechua para sus asegurados cusque{\~n}os \parencite{essalud2023}; en una regi{\'o}n
andina vecina, \textcite{cuaila2022} asoci{\'o} el dominio del quechua del personal con la calidad de
atenci{\'o}n percibida. En el plano tecnol{\'o}gico, el primer sistema de traducci{\'o}n neural para el
quechua se desarroll{\'o} con participaci{\'o}n de \emph{Hinantin Software} de Cusco
\parencite{ortega2020quechua}, y la variante Collao/Cusco (\texttt{quz}) fue incorporada reci{\'e}n en
la campa{\~n}a QUESPA 2025 \parencite{Ortega2025260}. A ello se suma el reto de la alfabetizaci{\'o}n en
zonas rurales andinas \parencite{inei2023,dispurse2023}. No se identific{\'o} ning{\'u}n antecedente
local que integre reconocimiento, traducci{\'o}n y s{\'i}ntesis de voz para la atenci{\'o}n en salud,
lo que confirma la pertinencia de esta investigaci{\'o}n.

\subsection{Estado del arte y brecha de investigaci{\'o}n}
A partir de una \textbf{revisi{\'o}n sistem{\'a}tica} (Scopus, IEEE Xplore y ACM Digital Library;
62 registros identificados, 58 tras depurar duplicados y 30 evaluados a texto completo), los
antecedentes de la secci{\'o}n anterior se sintetizan por componente en el Cuadro~\ref{tab:sintesis}
---sin repetir las citas ya detalladas---, lo que permite ubicar con precisi{\'o}n qu{\'e} se ha
resuelto y qu{\'e} sigue abierto.

\begin{table}[H]
\centering
\caption{S{\'i}ntesis del estado del arte por componente (elaboraci{\'o}n propia).}
\label{tab:sintesis}
\small
\begin{tabularx}{\linewidth}{@{}l X l@{}}
\toprule
\textbf{Componente} & \textbf{Avance / mejor resultado reportado} & \textbf{{\textquestiondown}Salud?} \\
\midrule
ASR (quechua)                  & WER 12.70\,\% (HMM/HTK) y 22.75\,\% con aumento de datos & No \\
MT (QU--ES)                    & Hasta $+$11\,\% chrF; BLEU 15.47 con modelos ajustados   & No \\
TTS (quechua)                  & MOS de hasta 4.9 (XTTS\,v2)                              & No \\
Traducci{\'o}n de voz (cascada) & 26.7 BLEU (QUESPA 2025, variante Collao)                & No \\
Recursos / corpus              & QuBERT, \emph{Huqariq}, API \emph{Siminchikkunarayku}    & Habilitador \\
\bottomrule
\end{tabularx}
\end{table}

La s{\'i}ntesis evidencia avances s{\'o}lidos por componente, pero ninguno en el \textbf{dominio
cl{\'i}nico} ni con la variante de Cusco (\texttt{quz}) integrada de extremo a extremo. La
\textbf{brecha} es, por tanto, la intersecci{\'o}n ---hoy vac{\'i}a--- entre las tecnolog{\'i}as del
habla para el quechua y la comunicaci{\'o}n m{\'e}dico-paciente: no existe un corpus cl{\'i}nico
Quechua--Espa{\~n}ol ni una evaluaci{\'o}n emp{\'i}rica de estos modelos en salud, y NLLB a{\'u}n no
cubre \texttt{quz}. Cerrar esa intersecci{\'o}n es el aporte de este proyecto.

\subsection{Fundamentos te{\'o}ricos}

\subsubsection{Tecnolog{\'i}as del habla y traducci{\'o}n autom{\'a}tica neuronal}
El \textbf{reconocimiento autom{\'a}tico de voz (ASR)} convierte una se{\~n}al de voz en texto; la
\textbf{traducci{\'o}n autom{\'a}tica (MT)}, hoy mayoritariamente \emph{neuronal} con arquitecturas
Transformer, traduce entre idiomas; y la \textbf{s{\'i}ntesis de voz (TTS)} genera voz a partir de
texto. En \textbf{lenguas de bajo recurso} ---como el quechua, de morfolog{\'i}a polisint{\'e}tica---
estas tareas se benefician de la transferencia desde modelos entrenados sobre grandes vol{\'u}menes de
datos multiling{\"u}es.

\subsubsection{Modelos multiling{\"u}es preentrenados (MMS, NLLB, Whisper)}
\textbf{MMS} \parencite{pratap2024mms} cubre el quechua de Cusco (\texttt{quz}) en reconocimiento y
s{\'i}ntesis; \textbf{NLLB} \parencite{costajussa2022nllb} ofrece traducci{\'o}n multiling{\"u}e que
incluye el quechua de Ayacucho (\texttt{quy}); y \textbf{Whisper} \parencite{radford2023whisper} se
emplea como referencia de ASR con cobertura marginal del quechua. Estos modelos permiten reutilizar
conocimiento ling{\"u}{\'i}stico general y especializarlo mediante ajuste.

\subsubsection{Ajuste fino y adaptaci{\'o}n dialectal (\texttt{quy}$\to$\texttt{quz})}
El \emph{ajuste fino} (\emph{fine-tuning}) adapta un modelo preentrenado a un dominio o variante
espec{\'i}ficos con relativamente pocos datos. El aporte t{\'e}cnico central de esta tesis es la
adaptaci{\'o}n del par de traducci{\'o}n del quechua de Ayacucho (\texttt{quy}) al de Cusco
(\texttt{quz}) sobre el vocabulario cl{\'i}nico, cerrando la brecha dialectal que hoy limita a NLLB.

\subsubsection{Arquitectura en cascada ASR$\to$MT$\to$TTS}
El encadenamiento de los tres componentes en una \textbf{arquitectura en cascada} habilita un
di{\'a}logo bidireccional por voz: la voz del hablante se transcribe (ASR), el texto se traduce (MT)
y la traducci{\'o}n se reproduce en la lengua del interlocutor (TTS). La modularidad facilita ajustar
y evaluar cada etapa por separado.

\subsubsection{M{\'e}tricas de evaluaci{\'o}n}
El desempe{\~n}o se mide con m{\'e}tricas estandarizadas: \textbf{WER} para el reconocimiento,
\textbf{BLEU} \parencite{papineni2002bleu} y \textbf{chrF} \parencite{popovic2015chrf} para la
traducci{\'o}n ---esta {\'u}ltima adecuada para lenguas morfol{\'o}gicamente ricas--- y \textbf{MOS}
para la naturalidad de la voz sintetizada. Su formulaci{\'o}n operativa se detalla en el
Cap{\'i}tulo~III.

\subsection{Definici{\'o}n de t{\'e}rminos b{\'a}sicos}
\begin{itemize}
  \item \textbf{Corpus:} conjunto estructurado de textos o grabaciones para entrenar o evaluar modelos.
  \item \textbf{Modelo preentrenado:} modelo entrenado sobre grandes vol{\'u}menes de datos, reutilizable mediante ajuste.
  \item \textbf{Lengua de bajo recurso:} idioma con escasez de datos digitales para el PLN.
  \item \textbf{Lengua polisint{\'e}tica:} lengua que concentra mucha informaci{\'o}n gramatical en una sola palabra, como el quechua.
  \item \textbf{Pipeline:} secuencia encadenada de etapas (voz $\to$ texto $\to$ traducci{\'o}n $\to$ voz).
\end{itemize}

% ##################################################################
% CAPÍTULO III: METODOLOGÍA DE LA INVESTIGACIÓN
% ##################################################################
\section{Metodolog{\'i}a de la investigaci{\'o}n}

\subsection{Enfoque, tipo y dise{\~n}o de investigaci{\'o}n}

\subsubsection{Enfoque de investigaci{\'o}n}
La investigaci{\'o}n adopta un enfoque \textbf{cuantitativo con complemento cualitativo} (mixto): las
m{\'e}tricas de desempe{\~n}o (WER, BLEU, chrF, MOS, latencia) aportan la evidencia cuantitativa, y la
percepci{\'o}n de utilidad y comprensi{\'o}n aporta el complemento cualitativo.

\subsubsection{Tipo de investigaci{\'o}n}
Es de tipo \textbf{aplicado} ---resuelve un problema concreto mediante un prototipo--- y de nivel
\textbf{explicativo}, pues contrasta el efecto del sistema frente a la ausencia de apoyo
tecnol{\'o}gico.

\subsubsection{Dise{\~n}o de investigaci{\'o}n}
El dise{\~n}o es \textbf{pre-experimental / cuasi-experimental}; la unidad de an{\'a}lisis es el
\emph{turno de traducci{\'o}n} dentro de una consulta simulada. La Figura~\ref{fig:arq} muestra el
flujo del sistema y el Algoritmo~\ref{alg:pipeline} detalla el procesamiento de un turno de
di{\'a}logo.

\begin{figure}[H]
\centering
\begin{tikzpicture}[
    node distance=0.6cm and 0.7cm,
    boxA/.style={rectangle, rounded corners, draw=g4, thick, fill=g3!18, text width=2.0cm, align=center, minimum height=1.2cm, font=\scriptsize},
    core/.style={rectangle, rounded corners, draw=g8, very thick, fill=g8!22, text width=2.0cm, align=center, minimum height=1.2cm, font=\scriptsize},
    boxB/.style={rectangle, rounded corners, draw=g11, thick, fill=g12!22, text width=2.0cm, align=center, minimum height=1.2cm, font=\scriptsize},
    io/.style={ellipse, draw=g5, thick, fill=white, text width=1.8cm, align=center, font=\scriptsize},
    bi/.style={{Stealth[length=2.5mm]}-{Stealth[length=2.5mm]}, thick, draw=g5}
]
\node[io] (pac) {Paciente\\(voz quechua)};
\node[boxA, right=of pac] (v1) {MMS\\STT / TTS\\(\texttt{quz})};
\node[core, right=of v1] (mt) {Traducci{\'o}n\\NLLB\\\texttt{quz}$\leftrightarrow$es};
\node[boxB, right=of mt] (v2) {MMS\\STT / TTS\\(es)};
\node[io, right=of v2] (med) {M{\'e}dico\\(voz espa{\~n}ol)};
\draw[bi] (pac) -- (v1);
\draw[bi] (v1) -- (mt);
\draw[bi] (mt) -- (v2);
\draw[bi] (v2) -- (med);
\end{tikzpicture}
\caption{Pipeline modular \textbf{bidireccional} por voz; la traducci{\'o}n NLLB adaptada al quechua de Cusco (\texttt{quz}) es el n{\'u}cleo del aporte. Elaboraci{\'o}n propia.}
\label{fig:arq}
\end{figure}

\begin{algorithm}[H]
\caption{Traducci{\'o}n bidireccional por voz de un turno de di{\'a}logo}
\label{alg:pipeline}
\begin{algorithmic}[1]
\Require Audio de entrada $a$; idioma origen $L_s$; idioma destino $L_t$
\Ensure Texto traducido $t_t$ y audio sintetizado $a_t$
\State $t_s \gets \textsc{STT}_{\mathrm{MMS}}(a,\,L_s)$
\State $t_t \gets \textsc{MT}_{\mathrm{NLLB}}(t_s,\,L_s \rightarrow L_t)$
\State Mostrar $t_t$ en pantalla
\State $a_t \gets \textsc{TTS}_{\mathrm{MMS}}(t_t,\,L_t)$
\State Reproducir $a_t$
\State \Return $(t_t,\,a_t)$
\end{algorithmic}
\end{algorithm}

\subsection{Variables y operacionalizaci{\'o}n}

\subsubsection{Identificaci{\'o}n de variables}
La \textbf{variable independiente} es el uso del sistema de traducci{\'o}n (con/sin sistema) y la
modalidad de salida (texto / texto+voz). Las \textbf{variables dependientes} son la precisi{\'o}n de
transcripci{\'o}n y traducci{\'o}n, la calidad de la voz y la comprensi{\'o}n/percepci{\'o}n. El
nivel de alfabetizaci{\'o}n act{\'u}a como \textbf{variable interviniente}.

\subsubsection{Matriz de operacionalizaci{\'o}n de variables}
El Cuadro~\ref{tab:variables} resume la operacionalizaci{\'o}n.

\begin{table}[H]
\centering
\caption{Operacionalizaci{\'o}n de variables.}
\label{tab:variables}
\small
\begin{tabularx}{\linewidth}{@{}l X l@{}}
\toprule
\textbf{Variable} & \textbf{Definici{\'o}n / dimensi{\'o}n} & \textbf{Indicador} \\
\midrule
\multirow{2}{*}{Independiente} & Uso del sistema de traducci{\'o}n IA & Con / sin sistema \\
                               & Modalidad de salida & Texto / texto+voz \\
\midrule
\multirow{3}{*}{Dependiente}   & Precisi{\'o}n de transcripci{\'o}n & WER \\
                               & Precisi{\'o}n de traducci{\'o}n & BLEU / chrF \\
                               & Comprensi{\'o}n y percepci{\'o}n & \% aciertos, Likert, MOS \\
\midrule
Interviniente                  & Nivel de alfabetizaci{\'o}n & Alto / bajo \\
\bottomrule
\end{tabularx}
\end{table}

\subsection{Unidades de an{\'a}lisis y selecci{\'o}n de consultas}

\subsubsection{Unidad de an{\'a}lisis}
La unidad de an{\'a}lisis es el \emph{turno de traducci{\'o}n} dentro de una consulta simulada
(una intervenci{\'o}n hablada que el sistema reconoce, traduce y reproduce). Los escenarios cubren
triaje, s{\'i}ntomas e indicaciones.

\subsubsection{Poblaci{\'o}n y muestra}
Participantes voluntarios (estudiantes de ciencias de la salud y/o personal asistencial
biling{\"u}e) que act{\'u}an roles en consultas simuladas; \emph{no} se usan pacientes ni datos
cl{\'i}nicos reales. Muestreo no probabil{\'i}stico por conveniencia (10--15 participantes; 30--50
frases-gui{\'o}n por escenario).

\subsection{T{\'e}cnicas e instrumentos de recolecci{\'o}n}
\begin{itemize}
  \item \textbf{Corpus:} lista de vocabulario y frases de triaje/s{\'i}ntomas/indicaciones, construida y revisada con personal de salud (fichas de validaci{\'o}n de contenido).
  \item \textbf{Registro de desempe{\~n}o:} \emph{logs} autom{\'a}ticos del sistema (WER, BLEU, latencia por turno).
  \item \textbf{Encuesta de percepci{\'o}n:} cuestionario tipo Likert (utilidad, confianza, comprensi{\'o}n) al personal y a los pacientes-actores.
  \item \textbf{Prueba de comprensi{\'o}n:} tarea de recuperaci{\'o}n de la indicaci{\'o}n (\% de aciertos) para contrastar H2 y H3.
\end{itemize}

\subsection{Procedimiento metodol{\'o}gico}
\begin{enumerate}
  \item Construcci{\'o}n y validaci{\'o}n del corpus m{\'e}dico Quechua--Espa{\~n}ol.
  \item Ajuste (\emph{fine-tuning}) de MMS (reconocimiento y s{\'i}ntesis) y, como aporte central, de NLLB para el par quechua de Cusco (\texttt{quz})--Espa{\~n}ol.
  \item Integraci{\'o}n del \emph{pipeline} y de la salida diferenciada por voz y texto.
  \item Ejecuci{\'o}n de consultas simuladas bajo las condiciones experimentales.
  \item Registro de m{\'e}tricas y aplicaci{\'o}n de instrumentos de percepci{\'o}n.
\end{enumerate}

\subsection{Protocolo de evaluaci{\'o}n}

\subsubsection{M{\'e}tricas de reconocimiento y traducci{\'o}n}
La precisi{\'o}n del reconocimiento se mide con la tasa de error de palabra (WER):
\begin{equation}
\mathrm{WER} = \frac{S + D + I}{N},
\end{equation}
\regecuacion{Tasa de error de palabra (WER)}
donde $S$, $D$ e $I$ son sustituciones, eliminaciones e inserciones sobre $N$ palabras de
referencia. La traducci{\'o}n se eval{\'u}a con BLEU \parencite{papineni2002bleu} y chrF
\parencite{popovic2015chrf}; esta {\'u}ltima, adecuada para lenguas morfol{\'o}gicamente ricas:
\begin{equation}
\mathrm{chrF}_{\beta} = (1 + \beta^{2})\,\frac{\mathrm{chrP}\cdot\mathrm{chrR}}{\beta^{2}\,\mathrm{chrP} + \mathrm{chrR}}.
\end{equation}
\regecuacion{Puntaje chrF ($F$-score de $n$-gramas de caracteres)}

\subsubsection{M{\'e}tricas de calidad de voz y comprensi{\'o}n}
La calidad de la voz sintetizada se valora con el \emph{Mean Opinion Score} (MOS). La comprensi{\'o}n
se mide con el porcentaje de aciertos en la recuperaci{\'o}n de la indicaci{\'o}n y con escalas Likert
de utilidad y confianza.

\subsubsection{M{\'e}tricas de eficiencia}
La eficiencia se cuantifica con la \emph{latencia por turno} (tiempo entre el fin de la voz de
entrada y el inicio de la voz de salida) y el consumo de recursos de c{\'o}mputo.

\subsection{M{\'e}todos de an{\'a}lisis de datos}
Estad{\'i}stica descriptiva para WER/BLEU/MOS/latencia y pruebas de comparaci{\'o}n de
medias/proporciones (t de Student o Wilcoxon seg{\'u}n normalidad), con un nivel de significancia
$\alpha = 0{,}05$, para contrastar las hip{\'o}tesis (modelo ajustado vs. sin ajustar; texto vs.
texto+voz).

\subsection{Aspectos {\'e}ticos y de gesti{\'o}n de datos}
El estudio usa {\'u}nicamente di{\'a}logos \textbf{simulados/actuados}, sin informaci{\'o}n cl{\'i}nica
real, evitando la necesidad de aprobaci{\'o}n de un comit{\'e} de {\'e}tica hospitalario para datos de
pacientes. Se obtiene consentimiento informado de los participantes voluntarios y se anonimizan las
grabaciones, en cumplimiento del atributo de {\'e}tica del curso. Los datos se almacenan de forma
segura y se usan solo con fines de investigaci{\'o}n.

% ##################################################################
% CAPÍTULO IV: ESTUDIO DE FACTIBILIDAD
% ##################################################################
\section{Estudio de factibilidad}

\subsection{Enfoque adoptado}
El proyecto es viable porque reutiliza modelos preentrenados de c{\'o}digo abierto (MMS, NLLB,
Whisper) mediante ajuste fino, lo que reduce dr{\'a}sticamente los requisitos de datos y c{\'o}mputo
frente a entrenar desde cero. La factibilidad se analiza en sus dimensiones t{\'e}cnica,
econ{\'o}mica y operativa.

\subsection{Factibilidad t{\'e}cnica}

\subsubsection{Hardware, c{\'o}mputo e inferencia}
El ajuste de modelos preentrenados puede realizarse en una estaci{\'o}n con GPU y complementarse con
cr{\'e}ditos de GPU en la nube para las fases intensivas. La inferencia del \emph{pipeline} integrado
es ejecutable en equipo de gama media, compatible con un piloto en el punto de atenci{\'o}n.

\subsubsection{Estrategia de datos}
El cuello de botella es el corpus cl{\'i}nico Quechua--Espa{\~n}ol, que se construir{\'a} de forma
acotada (vocabulario y frases de triaje, s{\'i}ntomas e indicaciones) y se complementar{\'a} con
recursos existentes (\emph{Huqariq}, QuBERT) y datos sint{\'e}ticos cuando sea pertinente.

\subsubsection{Herramientas y software}
Se emplear{\'a}n bibliotecas de c{\'o}digo abierto para habla y traducci{\'o}n (modelos MMS, NLLB y
Whisper de Meta/OpenAI y sus ecosistemas), herramientas de evaluaci{\'o}n (c{\'a}lculo de WER, BLEU y
chrF) y entornos reproducibles de \emph{fine-tuning}.

\subsection{Factibilidad econ{\'o}mica}
El presupuesto estimado (Cuadro~\ref{tab:presupuesto}) asciende a \textbf{S/ 7\,350.00}, dominado por
hardware con GPU y cr{\'e}ditos de c{\'o}mputo; es asumible en el marco de un trabajo de tesis y
puede financiarse con recursos propios y/o apoyo institucional.

\subsection{Factibilidad operativa}
La validaci{\'o}n con consultas simuladas y participantes voluntarios evita la dependencia de
aprobaciones cl{\'i}nicas complejas y permite ejecutar los experimentos dentro del semestre, en un
centro de salud de referencia en Cusco.

\subsection{Viabilidad y sostenibilidad}
El proceso (corpus + ajuste + evaluaci{\'o}n) es replicable en otras variantes del quechua o en otros
dominios, y el uso de modelos de c{\'o}digo abierto favorece la continuidad y la apropiaci{\'o}n local
de la tecnolog{\'i}a.

% ##################################################################
% CAPÍTULO V: ASPECTOS ADMINISTRATIVOS
% ##################################################################
\section{Aspectos administrativos}

\subsection{Cronograma de actividades}
El proyecto se ejecuta en \textbf{4 meses (octubre 2026 -- enero 2027)}, organizado por semanas. El
Cuadro~\ref{tab:cronograma} presenta el diagrama de Gantt, con las actividades agrupadas en tres fases
---\emph{administrativa}, de \emph{desarrollo} (ASR, MT y TTS) e \emph{investigaci{\'o}n y cierre}---;
los m{\'o}dulos ASR, MT y TTS se ejecutan de forma parcialmente paralela mediante el ajuste de modelos
preentrenados.

% ----- gris de la fase administrativa -----
\definecolor{adminG}{HTML}{8A98A6}
% \gbar{fila}{semana inicial}{semana final}{color}{etiqueta}
\newcommand{\gbar}[5]{%
  \fill[#4,rounded corners=2.2pt] ({#2-1},{-#1+1-0.20}) rectangle ({#3},{-#1+1+0.20});%
  \node[anchor=west,align=left,text width=4.9cm,font=\scriptsize,text=black!80]
       at (-8.2,{-#1+1}) {#5};}

\begin{table}[H]
\caption{Cronograma de actividades (octubre 2026 -- enero 2027), por semanas.}
\label{tab:cronograma}

% ----- leyenda de fases (alineada a la izquierda) -----
\noindent{\scriptsize
\begin{tikzpicture}[baseline]
  \fill[adminG,rounded corners=1.5pt](0,0) rectangle (0.32,0.26);
  \node[anchor=west] at (0.4,0.13) {Administrativo};
  \fill[g3,rounded corners=1.5pt](3.0,0) rectangle (3.32,0.26);
  \node[anchor=west] at (3.4,0.13) {Desarrollo (ASR, MT, TTS)};
  \fill[g10,rounded corners=1.5pt](7.6,0) rectangle (7.92,0.26);
  \node[anchor=west] at (8.0,0.13) {Investigaci{\'o}n y cierre};
\end{tikzpicture}}

\smallskip
\noindent\resizebox{\linewidth}{!}{%
\begin{tikzpicture}[x=0.62cm,y=0.82cm,font=\scriptsize]
  % --- banda de meses (navy) ---
  \fill[g5,rounded corners=1.5pt] (0,0.74) rectangle (16,1.6);
  \foreach \m/\pos in {Octubre/2,Noviembre/6,Diciembre/10,Enero/14}
     \node[white,font=\bfseries] at (\pos,1.17) {\m};
  % --- numeros de semana ---
  \foreach \w in {1,...,16}
     \node[font=\tiny,text=black!55] at ({\w-0.5},0.45) {\w};
  % --- separadores de mes ---
  \foreach \xg in {0,4,8,12,16}
     \draw[black!20,very thin] (\xg,0.74) -- (\xg,-11.6);
  % --- lineas horizontales entre filas ---
  \foreach \r in {0,...,12}
     \draw[black!10,very thin] (0,{-\r+0.5}) -- (16,{-\r+0.5});
  % --- barras (fila, semana inicial, final, color, etiqueta) ---
  \gbar{1}{1}{1}{adminG}{Aprobaci{\'o}n del plan de tesis}
  \gbar{2}{1}{4}{g1}{Revisi{\'o}n del estado del arte}
  \gbar{3}{2}{6}{g2}{Recolecci{\'o}n y preparaci{\'o}n del corpus (salud)}
  \gbar{4}{5}{8}{g3}{M{\'o}dulo ASR (voz $\to$ texto)}
  \gbar{5}{6}{9}{g4}{M{\'o}dulo MT (traducci{\'o}n)}
  \gbar{6}{7}{10}{g5}{M{\'o}dulo TTS (texto $\to$ voz)}
  \gbar{7}{10}{11}{g5}{Integraci{\'o}n del \emph{pipeline}}
  \gbar{8}{11}{13}{g8}{Evaluaci{\'o}n experimental (WER, BLEU, MOS)}
  \gbar{9}{12}{13}{g9}{Validaci{\'o}n m{\'e}dico-paciente (simulada)}
  \gbar{10}{13}{14}{g10}{An{\'a}lisis de resultados}
  \gbar{11}{13}{15}{g11}{Redacci{\'o}n del informe final}
  \gbar{12}{16}{16}{g12}{Sustentaci{\'o}n}
\end{tikzpicture}}
\end{table}

\subsection{Presupuesto y financiamiento}
El Cuadro~\ref{tab:presupuesto} resume los recursos estimados para el desarrollo del proyecto.

\begin{table}[H]
\centering
\caption{Recursos para el desarrollo del proyecto de tesis (soles).}
\label{tab:presupuesto}
\small
\begin{tabularx}{\linewidth}{@{}l X r@{}}
\toprule
\textbf{Tipo} & \textbf{Descripci{\'o}n} & \textbf{Costo (S/)} \\
\midrule
Hardware & Laptop/PC con GPU (entrenamiento y \emph{fine-tuning}) & 4\,500.00 \\
C{\'o}mputo & Cr{\'e}ditos de GPU en la nube (horas de ajuste) & 900.00 \\
Datos & Grabaci{\'o}n y validaci{\'o}n del corpus cl{\'i}nico (incentivos a hablantes) & 800.00 \\
Administrativos & Impresiones, inscripci{\'o}n y tr{\'a}mites de tesis & 600.00 \\
Movilidad & Traslados al centro de salud y a comunidades & 300.00 \\
Otros & Gastos no previstos & 250.00 \\
\midrule
\multicolumn{2}{@{}l}{\textbf{Total}} & \textbf{7\,350.00} \\
\bottomrule
\end{tabularx}
\end{table}

\subsection{Equipo de investigaci{\'o}n}
El Cuadro~\ref{tab:equipo} detalla el equipo responsable del proyecto.

\begin{table}[H]
\centering
\caption{Equipo de investigaci{\'o}n.}
\label{tab:equipo}
\small
\begin{tabularx}{\linewidth}{@{}l X@{}}
\toprule
\textbf{Rol} & \textbf{Responsable / funci{\'o}n} \\
\midrule
Tesista   & Alain Armando Condori Contreras ---dise{\~n}o, desarrollo y evaluaci{\'o}n del sistema---. \\
Asesor    & Mtro. Ing. Ra{\'u}l Huillca Huallparimachi ---orientaci{\'o}n metodol{\'o}gica y t{\'e}cnica---. \\
Colaboradores & Personal de salud biling{\"u}e (validaci{\'o}n del corpus) y participantes voluntarios (consultas simuladas). \\
\bottomrule
\end{tabularx}
\end{table}

% ================= BIBLIOGRAFÍA =================
\newpage
\printbibliography[heading=bibintoc,title={Referencias}]

% ================= ANEXOS =================
\newpage
\section*{Anexos}
\addcontentsline{toc}{section}{Anexos}

% ----- Anexo 1 -----
\subsection*{Anexo 1: Matriz de consistencia l{\'o}gica}
\addcontentsline{toc}{subsection}{Anexo 1: Matriz de consistencia l{\'o}gica}
Articula la coherencia interna entre el problema, los objetivos, las hip{\'o}tesis y las variables.

\begin{table}[H]
\centering
\caption{Matriz de consistencia l{\'o}gica.}
\label{tab:matriz}
\scriptsize
\renewcommand{\arraystretch}{1.3}
\begin{tabularx}{\linewidth}{@{}p{3.1cm} X X X@{}}
\toprule
\textbf{Problema} & \textbf{Objetivo} & \textbf{Hip{\'o}tesis} & \textbf{Variable/Indicador} \\
\midrule
General: {\textquestiondown}De qu{\'e} manera el sistema influye en la comunicaci{\'o}n m{\'e}dico-paciente? &
Desarrollar y evaluar el sistema de traducci{\'o}n bidireccional por voz. &
El sistema mejora significativamente la comunicaci{\'o}n m{\'e}dico-paciente. &
Comprensi{\'o}n (\% aciertos, Likert). \\
\midrule
E1: {\textquestiondown}Qu{\'e} corpus cl{\'i}nico se requiere? &
Construir y validar el corpus Quechua--Espa{\~n}ol. &
H1: el corpus validado aporta la base l{\'e}xica para el ajuste. &
Cobertura l{\'e}xica, validaci{\'o}n de contenido. \\
\midrule
E2: {\textquestiondown}El ajuste mejora la precisi{\'o}n? &
Adaptar NLLB al quechua de Cusco (\texttt{quz}) y MMS al dominio cl{\'i}nico. &
H2: el ajuste reduce el WER y mejora el BLEU/chrF. &
WER, BLEU/chrF. \\
\midrule
E3: {\textquestiondown}Qu{\'e} modalidad de salida es m{\'a}s comprensible? &
Integrar la salida por voz y texto bidireccional para ambos interlocutores. &
H3: texto+voz mejora la comprensi{\'o}n. &
\% aciertos, MOS. \\
\midrule
E4: {\textquestiondown}Qu{\'e} utilidad/confianza perciben los usuarios? &
Validar en consulta simulada y medir percepci{\'o}n. &
H4: mejora la comprensi{\'o}n mutua percibida. &
Escala Likert, tiempo. \\
\bottomrule
\end{tabularx}
\end{table}

% ----- Anexo 2 -----
\subsection*{Anexo 2: Matriz de consistencia metodol{\'o}gica}
\addcontentsline{toc}{subsection}{Anexo 2: Matriz de consistencia metodol{\'o}gica}
Conecta cada objetivo espec{\'i}fico con la variable que mide, la t{\'e}cnica e instrumento de
recolecci{\'o}n y el m{\'e}todo de an{\'a}lisis previsto.

\begin{table}[H]
\centering
\caption{Matriz de consistencia metodol{\'o}gica.}
\label{tab:matriz-met}
\scriptsize
\renewcommand{\arraystretch}{1.3}
\begin{tabularx}{\linewidth}{@{}p{3.1cm} X X X@{}}
\toprule
\textbf{Objetivo espec{\'i}fico} & \textbf{Variable / indicador} & \textbf{T{\'e}cnica e instrumento} & \textbf{M{\'e}todo de an{\'a}lisis} \\
\midrule
OE1. Construir y validar el corpus cl{\'i}nico QU--ES. &
Cobertura y validez del corpus (validaci{\'o}n de contenido). &
Fichas de validaci{\'o}n revisadas con personal de salud biling{\"u}e. &
Juicio de expertos; \% de concordancia. \\
\midrule
OE2. Ajustar (\emph{fine-tuning}) MMS y NLLB al dominio cl{\'i}nico. &
Precisi{\'o}n de reconocimiento y traducci{\'o}n (WER, BLEU, chrF). &
\emph{Logs} autom{\'a}ticos del sistema por turno. &
Estad{\'i}stica descriptiva; prueba t/Wilcoxon (ajustado vs. base). \\
\midrule
OE3. Integrar el \emph{pipeline} con salida por voz y texto. &
Calidad de voz y comprensi{\'o}n (MOS, \% aciertos). &
Prueba de comprensi{\'o}n; valoraci{\'o}n MOS. &
Comparaci{\'o}n texto vs. texto+voz. \\
\midrule
OE4. Validar en consultas simuladas y medir percepci{\'o}n. &
Utilidad, confianza y latencia (Likert, segundos). &
Cuestionario Likert; cron{\'o}metro/\emph{logs}. &
Estad{\'i}stica descriptiva; pruebas de comparaci{\'o}n. \\
\bottomrule
\end{tabularx}
\end{table}

% ----- Anexo 3 -----
\subsection*{Anexo 3: Ficha t{\'e}cnica de registro experimental}
\addcontentsline{toc}{subsection}{Anexo 3: Ficha t{\'e}cnica de registro experimental}
Plantilla para registrar cada \emph{turno de traducci{\'o}n} durante las consultas simuladas.

\begin{table}[H]
\centering
\caption{Ficha t{\'e}cnica de registro por turno (plantilla).}
\label{tab:ficha}
\small
\begin{tabularx}{\linewidth}{@{}p{4.6cm} X@{}}
\toprule
\textbf{Campo} & \textbf{Valor registrado} \\
\midrule
ID de turno / sesi{\'o}n            & \rule{5cm}{0.3pt} \\
Escenario (triaje/s{\'i}ntomas/indicaciones) & \rule{5cm}{0.3pt} \\
Direcci{\'o}n (\texttt{quz}$\to$es / es$\to$\texttt{quz}) & \rule{5cm}{0.3pt} \\
Transcripci{\'o}n ASR / WER         & \rule{5cm}{0.3pt} \\
Traducci{\'o}n MT / BLEU / chrF     & \rule{5cm}{0.3pt} \\
Calidad de voz (MOS)               & \rule{5cm}{0.3pt} \\
Latencia por turno (s)             & \rule{5cm}{0.3pt} \\
Observaciones                      & \rule{5cm}{0.3pt} \\
\bottomrule
\end{tabularx}
\end{table}

% ----- Anexo 4 -----
\subsection*{Anexo 4: Plantilla del banco de frases de consulta (evaluaci{\'o}n)}
\addcontentsline{toc}{subsection}{Anexo 4: Plantilla del banco de frases de consulta (evaluaci{\'o}n)}
Gui{\'o}n base de frases por escenario; el equivalente en quechua de Cusco se valida con hablantes
nativos antes de su uso.

\begin{table}[H]
\centering
\caption{Banco de frases de consulta simulada (ejemplos; \texttt{quz} por validar).}
\label{tab:banco}
\small
\begin{tabularx}{\linewidth}{@{}l X l@{}}
\toprule
\textbf{Categor{\'i}a} & \textbf{Enunciado (espa{\~n}ol)} & \textbf{Equivalente (\texttt{quz})} \\
\midrule
Triaje       & {\textquestiondown}Qu{\'e} lo trae a la consulta?          & [por validar] \\
Triaje       & {\textquestiondown}Hace cu{\'a}nto tiempo se siente mal?    & [por validar] \\
S{\'i}ntomas & {\textquestiondown}D{\'o}nde le duele?                      & [por validar] \\
S{\'i}ntomas & {\textquestiondown}Tiene fiebre o escalofr{\'i}os?          & [por validar] \\
Indicaciones & Tome esta pastilla cada 8 horas.                          & [por validar] \\
Indicaciones & Regrese en 3 d{\'i}as si no mejora.                        & [por validar] \\
\bottomrule
\end{tabularx}
\end{table}

% ----- Anexo 5 -----
\subsection*{Anexo 5: Cuestionario de percepci{\'o}n (escala Likert)}
\addcontentsline{toc}{subsection}{Anexo 5: Cuestionario de percepci{\'o}n (escala Likert)}
Escala de 1 (\emph{totalmente en desacuerdo}) a 5 (\emph{totalmente de acuerdo}); se aplica al
personal y a los pacientes-actores.

\begin{table}[H]
\centering
\caption{{\'I}tems del cuestionario de percepci{\'o}n.}
\label{tab:likert}
\small
\begin{tabularx}{\linewidth}{@{}c X c@{}}
\toprule
\textbf{N.\textordmasculine} & \textbf{{\'I}tem} & \textbf{1--5} \\
\midrule
1 & El sistema facilit{\'o} la comunicaci{\'o}n durante la consulta.     & \rule{1.6cm}{0.3pt} \\
2 & Conf{\'i}o en la traducci{\'o}n que entreg{\'o} el sistema.          & \rule{1.6cm}{0.3pt} \\
3 & Comprend{\'i} con claridad las indicaciones m{\'e}dicas.            & \rule{1.6cm}{0.3pt} \\
4 & La voz sintetizada fue clara y natural.                             & \rule{1.6cm}{0.3pt} \\
5 & El tiempo de respuesta del sistema fue adecuado.                    & \rule{1.6cm}{0.3pt} \\
6 & Usar{\'i}a el sistema en una consulta real.                         & \rule{1.6cm}{0.3pt} \\
\bottomrule
\end{tabularx}
\end{table}

% ----- Anexo 6 -----
\subsection*{Anexo 6: Modelo de consentimiento informado}
\addcontentsline{toc}{subsection}{Anexo 6: Modelo de consentimiento informado}
\noindent\textit{Yo, \rule{5cm}{0.3pt}, declaro haber sido informado(a) sobre los objetivos del
estudio ``Sistema bidireccional de traducci{\'o}n de voz Quechua--Espa{\~n}ol con IA preentrenada
para consultas m{\'e}dicas en Cusco''. Entiendo que participar{\'e} de forma voluntaria en consultas
\textbf{simuladas} ---sin datos cl{\'i}nicos reales---, que las grabaciones ser{\'a}n anonimizadas y
usadas solo con fines de investigaci{\'o}n, y que puedo retirarme en cualquier momento sin
perjuicio.}

\vspace{1.2cm}
\noindent\rule{6cm}{0.3pt}\hfill\rule{4cm}{0.3pt}\\
\noindent Firma del participante \hfill Fecha

\end{document}
